With the increasing demand for undergraduate theses and academic paper writing in universities, the standardization and efficiency of paper formatting have become important concerns for students and researchers. Although traditional word processing tools are easy to use, they have certain limitations in handling complex formulas, figure and table numbering, reference management, and template adaptation. LaTeX has strong academic typesetting capabilities, but its syntax rules, compilation environment, and error logs present a relatively high barrier for beginners. To address these problems, this thesis designs and implements an intelligent academic paper typesetting system based on JavaWeb and LaTeX, aiming to provide users with online, template-based, and intelligent support for academic writing and typesetting.

The system adopts a front-end and back-end separated architecture. The front end is built with Vue 3 and Vite to construct the user interaction interface, while the back end is developed using Spring Boot, MyBatis, and MySQL to implement business logic processing and data persistence. Centered on the academic writing process, the system mainly implements functions such as user registration and login, email verification codes, personal profile management, paper project management, online LaTeX editing, template management, paper compilation, PDF preview, file export, AI-assisted dialogue, and compilation error analysis. Users can create paper projects in the browser, select or import LaTeX templates, edit LaTeX source code online, and generate PDF documents through multiple compilation engines.

During system implementation, considering the large number of dependency files required by LaTeX templates and the difficulty of understanding compilation errors, the system supports importing ZIP template packages and preserves the template directory structure as much as possible during compilation, thereby reducing compilation failures caused by missing dependency files such as `.cls`, `.sty`, and `.bib` files. Meanwhile, the system introduces AI-assisted functions. When users encounter LaTeX syntax problems or compilation failures, the system can analyze error logs and provide possible causes and modification suggestions, thus lowering the difficulty of using LaTeX.

The test results show that the system can effectively complete tasks such as paper project management, template loading, online editing, paper compilation, PDF preview, and AI-based error analysis. It basically meets the needs of online academic paper typesetting and has certain practical value and application prospects.
